\documentclass[11pt,a4paper]{article}
\input{../../sablony/pripravy/priprava.sty}
\usepackage{tikz,pgfplots}
\pgfplotsset{compat=1.18}
\nastavPripravu{1. ročník}{Kvadratické rovnice a nerovnice}{127}{Smíšené kvadratické rovnice a nerovnice}
\begin{document}
\blok{Smíšené kvadratické rovnice a nerovnice}
Výběr metody z výrazu: odmocnění, vzorec, rozklad, znaménková tabulka nebo graf.
\blok{Metoda a návaznost na graf}
\begin{itemize}
\item V každé úloze zdůvodni výběr metody a zapisuj množinu $K$.
\item Úlohy rozděl podle individuálních slabin, neřeš celou sbírku frontálně.
\end{itemize}
\blok{Společně řešené příklady}
\begin{prikladReseniBox}[title={Příklad 1}]
$x^{2} - 5x + 6=0$
\tcblower
\[
\begin{aligned}
D&=25-4\cdot(1)\cdot(6)=1\\[2pt]
x_1&=2\\[2pt]
x_2&=3\\[2pt]
K&=\{2;3\}
\end{aligned}
\]
\end{prikladReseniBox}
\begin{prikladReseniBox}[title={Příklad 2}]
$x^{2} - 5x + 6<0$
\tcblower
\[
\begin{aligned}
f(x)&=x^{2} - 5x + 6\\[2pt]
f(x)=0&\ \Rightarrow\ K_0=\{2;3\}\\[2pt]
K&=(2;3)
\end{aligned}
\]
\end{prikladReseniBox}
\begin{prikladReseniBox}[title={Příklad 3}]
$\left|x^{2} - 4\right|=3$
\tcblower
\[
\begin{aligned}
|f(x)|&=3\\[2pt]
f(x)&=3\quad\vee\quad f(x)=-3\\[2pt]
K_1&=\{- \sqrt{7};\sqrt{7}\}\\[2pt]
K_2&=\{-1;1\}\\[2pt]
K&=\{- \sqrt{7};-1;1;\sqrt{7}\}
\end{aligned}
\]
\end{prikladReseniBox}
\blok{Samostatné procvičování}
\begin{zadaniReseniBox}[title={Samostatná práce 1}]
$x^{4} - 5x^{2} + 4=0$
\tcblower
\[
\begin{aligned}
t&=x^2\ge0\\[2pt]
0&=t^{2} - 5 \cdot t + 4\\[2pt]
t&\in\{1;4\}\\[2pt]
K&=\{-2;-1;1;2\}
\end{aligned}
\]
\end{zadaniReseniBox}
\begin{zadaniReseniBox}[title={Samostatná práce 2}]
$2x^{2} - 3x - 2=0$
\tcblower
\[
\begin{aligned}
D&=9-4\cdot(2)\cdot(-2)=25\\[2pt]
x_1&=- \frac{1}{2}\\[2pt]
x_2&=2\\[2pt]
K&=\{- \frac{1}{2};2\}
\end{aligned}
\]
\end{zadaniReseniBox}
\blok{Volitelná rezerva}
\poznamka{Studenti si vymyslí obdobný příklad se změnou jednoho koeficientu a vysvětlí, jak se změní počet kořenů či řešení.}
\blok{Rozšiřující příklady ze sbírky}
\poznamka{Připomenutí nulového součinu před souhrnným procvičováním.}

\begin{zadaniReseniBox}[title={Rozšíření 1}]
Řešte v $\mathbb R$: $- \left(4x^{2} - 9\right) \cdot \left(x^{2} - 6x + 9\right)=0$.
\tcblower
\[
\begin{aligned}- \left(x - 3\right)^{2} \cdot \left(2x - 3\right) \cdot \left(2x + 3\right)&=0\\K&=\{- \frac{3}{2};\frac{3}{2};3\}\end{aligned}
\]
\end{zadaniReseniBox}

\blok{Volitelná výzva -- substituce}
\begin{zadaniReseniBox}[title={Rozšíření 2}]
$x^4-5x^2+4=0$.
\tcblower
\[\begin{aligned}
t&=x^2\ge0\\t^2-5t+4&=0\\t&=1\ \text{nebo}\ 4\\K&=\{-2;-1;1;2\}
\end{aligned}\]
\end{zadaniReseniBox}
\blok{Domácí úkol}
\begin{zadaniReseniBox}[title={Domácí úkol}]
\begin{enumerate}[label=\alph*),itemsep=2pt]
\item $x^{2} - 16=0$
\item $x^{2} + 2x + 5=0$
\item $x^{2} - 9<0$
\item $x^{2} - 4 \cdot \left|{x}\right| + 3=0$
\item $x^{4} - 10x^{2} + 9=0$
\item $\dfrac{\left(x - 2\right) \cdot \left(x + 2\right)}{x - 1}>0$
\end{enumerate}
\tcblower
\begin{enumerate}[label=\alph*),itemsep=2pt]
\item
\[\begin{aligned}
x^2&=16\\|x|&=4\\[2pt]
K&=\{-4;4\}
\end{aligned}\]
\item
\[\begin{aligned}
D&=4-4\cdot(1)\cdot(5)=-16\\[2pt]
D&<0\ \Rightarrow\ \text{žádný reálný kořen}\\[2pt]
K&=\varnothing
\end{aligned}\]
\item
\[\begin{aligned}
f(x)&=x^{2} - 9\\[2pt]
f(x)=0&\ \Rightarrow\ K_0=\{-3;3\}\\[2pt]
K&=(-3;3)
\end{aligned}\]
\item
\[\begin{aligned}
t&=|x|\ge0\\[2pt]
t^{2} - 4 \cdot t + 3&=0\quad\text{(řešíme pro }t\text{)}\\[2pt]
t&\in \{1;3\}\\[2pt]
K&=\{-3;-1;1;3\}
\end{aligned}\]
\item
\[\begin{aligned}
t&=x^2\ge0\\[2pt]
0&=t^{2} - 10 \cdot t + 9\\[2pt]
t&\in\{1;9\}\\[2pt]
K&=\{-3;-1;1;3\}
\end{aligned}\]
\item
\[\begin{aligned}
x&\ne 1\quad\text{(podmínka)}\\[2pt]
\text{Kritické body: }&-2,\ 1,\ 2\\[2pt]
K&=(-2;1)\cup(2;+\infty)
\end{aligned}\]
\end{enumerate}
\end{zadaniReseniBox}
\poznamka{Orientační čas: úvod 5 min, společné příklady 18 min, samostatně 15 min, kontrola a shrnutí 7 min. Rezerva je dobrovolná. Číselné parametry modelových aplikací nejsou skutečně naměřená data.}
\end{document}
